% === Begin Label Data ===
% ID: 876
% 难度星级: 
% 标签: 
% 备注: 
% 组卷引用次数: 0
% === End  Label Data ===

\begin{problem}{2024}{G}{北京卷}{21}{数列}
已知集合 $ M=\{(i,j,k,w)\mid i\in\{1,2\},j\in\{3,4\},k\in\{5,6\},w\in\{7,8\} \text{且} i+j+k+w \text{为偶数}\} $. 给定数列 $ A:a_1,a_2,\cdots,a_8 $ 和序列 $ \Omega:T_1,T_2,\cdots,T_s $，其中 $ T_t=(i_t,j_t,k_t,w_t)\in M(t=1,2,\cdots,s) $，对数列 $ A $ 进行如下变换：将 $ A $ 的第 $ i_t,j_t,k_t,w_t $ 项均加 $ 1 $，其余项不变，得到的数列记作 $ T_t(A) $；将 $ T_t(A) $ 的第 $ i_2,j_2,k_2,w_2 $ 项均加 $ 1 $，其余项不变，得到的数列记作 $ T_2T_1(A) $；……以此类推. 得到数列 $ T_s\cdots T_2T_1(A) $，简记为 $ \Omega(A) $.

（1）给定数列 $ A:1,3,2,4,6,3,1,9 $ 和序列 $ \Omega:(1,3,5,7),(2,4,6,8),(1,3,5,7) $，写出 $ \Omega(A) $；

（2）是否存在序列 $ \Omega $，使得 $ \Omega(A):a_1+2,a_2+6,a_3+4,a_4+2,a_5+8,a_6+2,a_7+4,a_8+4 $？若存在，写出一个 $ \Omega $；若不存在，请说明理由；

（3）若数列 $ A $ 的各项均为正整数，且 $ a_1+a_3+a_5+a_7 $ 为偶数. 求证：“存在序列 $ \Omega $，使得 $ \Omega(A) $ 的各项都相等”的充要条件为“$ a_1+a_2=a_3+a_4=a_5+a_6=a_7+a_8 $”.
\end{problem}

\begin{answer}
（1）$ \Omega(A):2,3,3,4,7,3,2,9 $
\end{answer}

\begin{solutions}
（1）$ T_1(A):2,3,3,4,7,3,2,10 $
$ T_2T_1(A):2,4,3,5,7,4,2,10 $
$ T_3T_2T_1(A):3,4,4,5,8,4,3,10 $
即 $ \Omega(A):3,4,4,5,8,4,3,10 $

（2）不存在. 理由如下：
若存在 $ \Omega $，则 $ a_1,a_2 $ 与 $ a_3,a_4 $ 增加值之和应该相等.
注意到 $ a_1,a_2 $ 一共增加了 $ 8 $，而 $ a_3,a_4 $ 一共增加了 $ 6 $，从而不存在符合题意的 $ \Omega $.

（3）将 $ A, T_1(A), T_2T_1(A), \cdots, T_s\cdots T_2T_1(A) $ 写成 $ s+1 $ 行 $ 8 $ 列的数表 $ N $.
用 $ N(x,y) $ 表示第 $ x $ 行 $ y $ 列的数，其中 $ x,y\in\mathbb{N}^*,x\leqslant s,y\leqslant 8 $.
必要性：根据 $ T $ 变换的定义，每一行第 $ 1,2 $ 列、第 $ 3,4 $ 列、第 $ 5,6 $ 列、第 $ 7,8 $ 列之和均比前一列对应和多 $ 1 $. 而 $ \Omega(A) $ 的各项均相等，设最后一行第 $ 1,2 $ 列数之和均为 $ m $，因此最后一行所有数相同，因此第 $ 1 $ 行第 $ 1,2 $ 列、第 $ 3,4 $ 列、第 $ 5,6 $ 列、第 $ 7,8 $ 列之和均为 $ m-s $，必要性得证.
充分性：若 $ a_1+a_2=a_3+a_4=a_5+a_6=a_7+a_8=m $，则根据 $ T $ 变换的定义，有
$$ N(k,1)+N(k,2)=N(k,3)+N(k,4) $$
$$ =N(k,5)+N(k,6)=N(k,7)+N(k,8)=m+(k-1). $$
且 $ N(k,1)+N(k,3)+N(k,5)+N(k,7) $ 始终为偶数.
接下来证明可以通过每次选择在 $ N(k,1),N(k,3),N(k,5),N(k,7) $ 选择其中 $ 2 $ 个使其增加 $ 1 $ 的方式，最终使得 $ \Omega(A) $ 的各项均相等，即 $ N(k,1)=N(k,3)=N(k,5)=N(k,7) $.
则考虑 $ N(k,1),N(k,3),N(k,5),N(k,7) $ 中最大值与最小值之差 $ d $.
若 $ d\neq 0 $，则按 $ N(k,1),N(k,3),N(k,5),N(k,7) $ 中的不同取值个数分类.
情形一：$ N(k,1),N(k,3),N(k,5),N(k,7) $ 有 $ 3 $ 个或 $ 4 $ 个不同取值.
此时将最小的两个分别加 $ 1 $，则 $ d $ 会减小 $ 1 $.
情形二：$ N(k,1),N(k,3),N(k,5),N(k,7) $ 有 $ 2 $ 个不同取值.
设取值为 $ m,n,m,n $，此时有三类：
(1) 取值为 $ m,m,n,n $，此时将其变为 $ m,m,n+1,n+1 $，则 $ d $ 会减小 $ 1 $.
(2) 取值为 $ m,m,n,n $，此时 $ 3m+n $ 为偶数，于是 $ m-n $ 为偶数，从而 $ m\geqslant n+2 $，进而将其变为 $ m+1,m,n+1\to m+1,m+1,m,n+2 $，则 $ d\to(m+1)-(n+2)=m-n-1 $，$ d $ 会减小 $ 1 $.
(3) 取值为 $ m,n,n,n $，此时 $ m+3n $ 为偶数，于是 $ m-n $ 为偶数，从而 $ m\geqslant n+2 $，进而将其变为 $ m,n+1,n\to m,n+2,n+1,n+1 $，则 $ d\to m-(n+1)=m-n-1 $，$ d $ 会减小 $ 1 $.
综上所述，一定会减小到 $ 0 $，证明完成.
\end{solutions}